\section{The Propers: Themes and Movement}
\zlabel{triptych:concise:themes:start}
God's gifts make a merciful life possible and sustain it through a journey that remains dangerous. The forgiven servant must become merciful; the protected family must persevere without supposing that every wound proves abandonment. The Gospel's guilty debtor and the Offertory's upright Job give these two readings their different centers.

\subsection{The family receives before it acts}
Mordecai invokes the creator while his people face destruction. The Introit's sovereignty belongs to God, not to the persecutor who appears to control their future. Its Psalm verse names a walk in the law; the Collect asks that continuing mercy protect the family and free it for good works. Obedience is a life made possible by help. Augustine notices the same dependence in Psalm~118: the command to walk is accompanied by a prayer that God direct the walk.\footnote{Esth.~3--4;13:8--18; Augustine, Ps.~118, Aleph, especially \S\S3,5--7.}

Chrysostom's account of the Gospel begins from God's prior gifts: creation, redemption, baptismal pardon and the Spirit. The king's cancellation of an impossible debt makes the servant's release real. Augustine, in turn, identifies the hearer as both a debtor to God and a creditor of another person. Mercy is therefore neither a possession to hoard nor a reward for having first made oneself independent. These Fathers explain different parts of the same demand: remember what was given, and recognize the dependence that continues whenever you pray.\footnote{Chrysostom, \work{Matthew} LXI, \S\S1--4; Augustine, Sermon~83 \S\S1--4.}

The Gradual gives this family an abiding refuge. Generations pass; mountains were made; God is. For Augustine, the changing person begins to take refuge in God, while God does not begin to exist or become faithful at that moment. The first reading finds here room for repentance after failure. The second finds an answer to the fear that a changing life means a changing or defeated protector.\footnote{Augustine, Ps.~89 \S\S2--4.}

\subsection{Freedom meets a brother and a battle}
The king gives more than the servant requests: pardon instead of delay. Outside, the fellow servant repeats the plea for patience and receives less than he asks: prison. The first man's new freedom has become an instrument of cruelty. Chrysostom makes the repeated words decisive; Augustine insists that pardon must reach the heart. Their warning falls on a deliberate refusal of mercy, not on the mere presence of grief or anger in an injured person.

Paul's armor exposes the wrong direction of that violence. The struggle is not against flesh and blood. Truth, justice, faith and the Gospel of peace prepare a person to resist evil without treating a neighbor as a being beyond charity. Chrysostom recalls the unmerciful servant when explaining the common heavenly Master immediately before the armor passage. Earthly rank cannot exempt anyone from that Lord's judgment.\footnote{Chrysostom, \work{Ephesians} XXII, on 6:9--13; Augustine, Sermon~83 \S8.}

The Alleluia's departure from Egypt adds another question: what follows liberation? Augustine reads the sea as baptismal deliverance and recalls Paul's warning that many later fell in the wilderness. Chrysostom reads the girded traveler through the Passover and the Church's pilgrimage in Christ. The delivered still need armor. The Gospel gives that danger its moral edge: perseverance includes mercy, not merely survival of other people's hostility.

\newpage
\subsection{The wounded still belong to their helper}
Job is called upright before his possessions and children are lost and his body is struck. His suffering cannot therefore be read as the Gospel servant's punishment under another name. Gregory distinguishes the enemy's wish to injure from the power permitted him: Satan does not possess an independent sovereignty. The successive permissions are bounded, and Job can remain under God's guardianship while really suffering.\footnote{Gregory, \work{Moralia} II.10.16--11.19; III.5.6--7.}

This distinction gives the readings different uses for the same Offertory. In the reading of received mercy, Job prevents an accusation: not every afflicted person is an unmerciful debtor. In the reading of protected pilgrimage, Job answers the central question: divine protection need not mean immunity from wounds. Neither answer cancels the Collect's plea for concrete rescue. The family may ask for adversity to cease, and its members may work to protect one another.

The Secret then seeks reconciliation and restored salvation through the offering God himself has willed. Human endurance is not the saving remedy, and merciful conduct does not purchase God's first gift. Schuster reads Job toward the innocent, suffering Christ and the sacrifice's restoration. The fulfillment in Christ deepens Job's significance without turning his literal affliction into proof of personal guilt. The Trinity Preface confesses the one God in distinct Persons and equal majesty: mercy, justice and saving power are not competing divine interests.\footnote{Schuster, \work{Sacramentary} III, pp.~178,181; Augustine, \work{On the Trinity} III, preface \S3.}

\subsection{The food has a destination and an inward fruit}
At Communion the people still ask for judgment against persecution. Augustine hears the whole Christ longing for salvation; Bellarmine hears the desire for everlasting peace. Mercy does not require naming injustice good, and hope does not require pretending the desired help has already become visible. Judgment entrusted to God differs from the forgiven servant's seizure of another's throat.\footnote{Augustine, Ps.~118, Caph, on vv.~81--88; Bellarmine, Ps.~118:81--88.}

The Postcommunion names the gift as nourishment of immortality and asks that the mind follow what the mouth received. For the first reading, the decisive fruit is a heart made merciful. For the second, the food sustains the passage toward the rest for which it is given. The sacramental gift is real in both; reception opens into a life and does not make that life unnecessary. Schuster draws the prayer into mind, heart and will.\footnote{Schuster, III, p.~181; Chrysostom, \work{Ephesians} XXIII.}

The conditional prayers extend this movement. The mission Collect asks workers and confident speech; its Secret seeks worship among all nations through Christ's redemption; its Postcommunion asks growth in true faith. Salvation goes outward and must continue to bear fruit. Where Luke is commemorated, his Collect names cross-bearing service, his Secret joins healing to free service and glory, and his Postcommunion asks sanctification and safety through the altar's gift. The saint's costly service and the family's hope belong together. The Mass ends with neither the debtor's closed fist nor the pilgrim's abandonment, but with grace received for a life still to be lived.
\zlabel{triptych:concise:themes:end}
\clearpage
